% !TEX program = pdflatex
\documentclass[10pt,aspectratio=169]{beamer}

\usepackage[utf8]{inputenc}
\usepackage[T1]{fontenc}
\usepackage[english]{babel}
\usepackage{lmodern}
\usepackage{microtype}
\usepackage{amsmath,amssymb,bm,mathtools}
\usepackage{booktabs,array}

\usetheme[numbering=fraction,progressbar=frametitle]{metropolis}
\setbeamertemplate{frame numbering}[fraction]
\setbeamertemplate{navigation symbols}{}

\definecolor{CourseBlue}{HTML}{1A3C68}
\definecolor{CourseGreen}{HTML}{216E39}
\definecolor{CourseGold}{HTML}{B8860B}
\definecolor{CourseRed}{HTML}{9A2C2C}
\definecolor{SoftBlue}{HTML}{F2F6FA}
\definecolor{SoftGreen}{HTML}{F1F8F3}
\definecolor{SoftRed}{HTML}{FFF3F3}

\setbeamercolor{normal text}{fg=black,bg=white}
\setbeamercolor{structure}{fg=CourseBlue}
\setbeamercolor{title}{fg=CourseBlue}
\setbeamerfont{title}{size=\Large,series=\bfseries}
\setbeamercolor{frametitle}{bg=CourseBlue,fg=white}
\setbeamerfont{frametitle}{series=\bfseries,size=\large}
\setbeamertemplate{frametitle}{
  \nointerlineskip
  \begin{beamercolorbox}[wd=\paperwidth,ht=3.2ex,dp=1.4ex,leftskip=1em,rightskip=1em]{frametitle}
    \usebeamerfont{frametitle}\insertframetitle
  \end{beamercolorbox}
}
\setbeamercolor{block title}{bg=CourseBlue,fg=white}
\setbeamercolor{block body}{bg=SoftBlue,fg=black}
\setbeamercolor{block title example}{bg=CourseGreen,fg=white}
\setbeamercolor{block body example}{bg=SoftGreen,fg=black}
\setbeamercolor{block title alerted}{bg=CourseRed,fg=white}
\setbeamercolor{block body alerted}{bg=SoftRed,fg=black}
\setbeamertemplate{blocks}[rounded][shadow=false]

\newenvironment{modelbox}{\begin{beamercolorbox}[rounded=true,sep=0.7em,wd=\textwidth]{block body}
  \textbf{\color{CourseBlue}MODEL ASSUMPTION}\quad}{\end{beamercolorbox}}
\newenvironment{marketbox}{\begin{beamercolorbox}[rounded=true,sep=0.7em,wd=\textwidth]{block body example}
  \textbf{\color{CourseGreen}MARKET ASSUMPTION}\quad}{\end{beamercolorbox}}
\newenvironment{simplificationbox}{\begin{beamercolorbox}[rounded=true,sep=0.7em,wd=\textwidth]{block body alerted}
  \textbf{\color{CourseRed}CHECK THE SIGN}\quad}{\end{beamercolorbox}}
\newenvironment{pausebox}{\begin{beamercolorbox}[rounded=true,sep=0.8em,wd=\textwidth]{block body}
  \textbf{\color{CourseGold}YOUR TURN}\quad}{\end{beamercolorbox}}

\newcommand{\dd}{\,\mathrm d}
\newcommand{\E}{\mathbb E}
\newcommand{\Pp}{\mathbb P}
\newcommand{\Q}{\mathbb Q}
\newcommand{\R}{\mathbb R}
\newcommand{\cF}{\mathcal F}
\newcommand{\1}{\mathbf 1}
\newcommand{\Nlaw}{\mathcal N}
\DeclareMathOperator{\Var}{Var}
\DeclareMathOperator{\Cov}{Cov}

\AtBeginSection[]{
  \begin{frame}[plain]
    \vfill
    {\Large\bfseries\color{CourseBlue}\insertsectionhead\par}
    \medskip
    {\color{CourseBlue}\rule{0.92\textwidth}{0.8pt}}
    \vfill
  \end{frame}
}

\title{Girsanov and Risk-Neutral Pricing}
\subtitle{Changing probabilities, moving drifts, and pricing by expectation}
\author{}
\date{Lecture 6}

\begin{document}

\begin{frame}[plain]
  \vfill
  {\LARGE\bfseries\color{CourseBlue}Girsanov and Risk-Neutral Pricing\par}
  \medskip
  {\large Changing probabilities, moving drifts, and pricing by expectation\par}
  \medskip
  {\color{CourseGold}\rule{0.94\textwidth}{0.8pt}}
  \medskip
  {Lecture 6\par}
  \vfill
\end{frame}

\begin{frame}{One question drives the whole lecture}
\begin{block}{The question}
Can we change probabilities so that the same risky asset has drift \(r\)
instead of drift \(\mu\)?
\end{block}
\begin{enumerate}
  \item Reweight scenarios and derive Bayes' rule.
  \item See Gaussian tilting move a mean.
  \item Extend the idea to Brownian paths.
  \item Build the stochastic exponential and state Girsanov.
  \item Apply the theorem to discounted stock prices.
  \item Recover pricing by conditional expectation.
\end{enumerate}
\end{frame}

\section{Reweighting Scenarios and Bayes' Rule}

\begin{frame}{A Gaussian tail is invisible to direct Monte Carlo}
Let \(X\sim\Nlaw(0,1)\) under \(\Pp\). We want
\[
p=\Pp(X>5)=1-\Phi(5)\approx 2.87\times10^{-7}.
\]
With \(n=10^6\) independent draws, direct Monte Carlo has
\[
\widehat p_{\mathrm{MC}}=\frac1n\sum_{i=1}^n\1_{\{X_i>5\}},
\qquad \E^\Pp[\text{number of hits}]=np\approx0.29.
\]
\begin{alertblock}{Too few tail observations}
The chance of seeing no hit is \((1-p)^n\approx75\%\); then the estimate is zero.
\end{alertblock}
\begin{exampleblock}{Change the sampling probabilities}
Give more weight to the tail: under a new law \(\Q\), let
\(X\sim\Nlaw(5,1)\), so \(\Q(X>5)=1/2\). We will correct for the change
of probabilities after sampling.
\end{exampleblock}
\end{frame}

\begin{frame}{On a finite space a density is a probability ratio}
Let \(\Omega=\{\omega_1,\ldots,\omega_m\}\), with positive probabilities
under \(\Pp\) and \(\Q\). Define
\[
Z(\omega_i)=\frac{\Q(\omega_i)}{\Pp(\omega_i)}.
\]
Then
\[
Z>0,\qquad \E^\Pp[Z]=1,\qquad
\boxed{\E^\Q[Y]=\E^\Pp[ZY].}
\]
\begin{block}{Interpretation}
\(Z>1\) upweights a scenario; \(Z<1\) downweights it.
\end{block}
\end{frame}

\begin{frame}{A three-state regime}
\[
\begin{array}{c|ccc}
 & \omega_1 & \omega_2 & \omega_3\\ \hline
\Pp & 1/2 & 1/3 & 1/6\\
\Q  & 1/4 & 1/2 & 1/4
\end{array}
\]
\begin{pausebox}
Compute \(Z\), identify which states are upweighted, and check
\(\E^\Pp[Z]=1\).
\end{pausebox}
\pause
\[
\boxed{Z(\omega_1)=\tfrac12,\quad
Z(\omega_2)=\tfrac32,\quad Z(\omega_3)=\tfrac32.}
\]
\end{frame}

\begin{frame}{Use the density rather than recomputing every probability}
In the previous model, let \(Y=(-2,1,4)\).
\begin{pausebox}
Compute \(\E^\Q[Y]\) directly, then as \(\E^\Pp[ZY]\).
\end{pausebox}
\pause
\[
\E^\Q[Y]=-\frac12+\frac12+1=1,
\]
\[
\E^\Pp[ZY]
=\frac12\frac12(-2)+\frac13\frac32(1)+\frac16\frac32(4)=1.
\]
\end{frame}

\begin{frame}{The general construction needs positivity and unit mean}
Take \(Z\ge0\) with \(\E^\Pp[Z]=1\), and define
\[
\Q(A)=\E^\Pp[Z\1_A].
\]
Then \(Z=\dd\Q/\dd\Pp\) and
\[
\boxed{\E^\Q[Y]=\E^\Pp[ZY].}
\]
\begin{alertblock}{Equivalence}
If \(Z>0\) almost surely, \(\Pp\) and \(\Q\) have the same null events.
No possible scenario is created or deleted.
\end{alertblock}
\end{frame}

\begin{frame}{Conditional expectations need conditional normalization}
For \(\mathcal G\subseteq\cF\),
\[
\boxed{
\E^\Q[Y\mid\mathcal G]
=\frac{\E^\Pp[ZY\mid\mathcal G]}
{\E^\Pp[Z\mid\mathcal G]}.}
\]
\begin{block}{Why the denominator appears}
It renormalizes the weights after the information in \(\mathcal G\) has
been revealed.
\end{block}
\end{frame}

\begin{frame}{Prove Bayes' rule from its defining test}
Set
\[
X=\frac{\E^\Pp[ZY\mid\mathcal G]}
{\E^\Pp[Z\mid\mathcal G]}.
\]
\begin{pausebox}
For \(A\in\mathcal G\), show that
\(\E^\Q[\1_AX]=\E^\Q[\1_AY]\).
\end{pausebox}
\pause
\[
\E^\Pp[\1_AX\E^\Pp[Z\mid\mathcal G]]
=\E^\Pp[\1_A\E^\Pp[ZY\mid\mathcal G]]
=\E^\Pp[\1_AZY].
\]
\end{frame}

\begin{frame}{Known information is unchanged by reweighting}
If \(Y\) is \(\mathcal G\)-measurable, then
\[
\E^\Pp[ZY\mid\mathcal G]
=Y\E^\Pp[Z\mid\mathcal G],
\]
so
\[
\boxed{\E^\Q[Y\mid\mathcal G]=Y.}
\]
\begin{pausebox}
Recover the unconditional identity by taking the trivial sigma-field.
\end{pausebox}
\end{frame}

\section{Gaussian Tilting Moves a Mean}

\begin{frame}{The one-dimensional prototype fixes the sign convention}
Let \(X\sim\Nlaw(0,1)\) under \(\Pp\) and
\[
Z=\exp\!\left(-\theta X-\frac12\theta^2\right).
\]
The moment-generating function (MGF) is
\(M_X(u)=\E^\Pp[e^{uX}]=e^{u^2/2}\). Hence
\[
\E^\Pp[Z]=e^{-\theta^2/2}M_X(-\theta)=1.
\]
\begin{pausebox}
For \(\theta>0\), does \(Z\) favor positive or negative \(X\)?
\end{pausebox}
\pause
It favors negative values because the weight contains \(e^{-\theta X}\).
\end{frame}

\begin{frame}{Identifies the new law}
\[
e^{-\theta x-\theta^2/2}\frac{e^{-x^2/2}}{\sqrt{2\pi}}
=\frac1{\sqrt{2\pi}}e^{-(x+\theta)^2/2}.
\]
Therefore
\[
\boxed{X\sim\Nlaw(-\theta,1)\text{ under }\Q.}
\]
\begin{exampleblock}{Prototype of Girsanov}
The mean changes; the variance stays equal to one.
\end{exampleblock}
\end{frame}

\begin{frame}{Reverse-engineer a target Gaussian mean}
Suppose \(X\sim\Nlaw(0,\sigma^2)\) under \(\Pp\).
\begin{pausebox}
Find \(Z=e^{aX+b}\) that makes \(X\sim\Nlaw(m,\sigma^2)\) under \(\Q\).
\end{pausebox}
\pause
\[
\boxed{
a=\frac{m}{\sigma^2},
\qquad
b=-\frac{m^2}{2\sigma^2}.}
\]
The term \(b\) enforces \(\E^\Pp[Z]=1\).
\end{frame}

\begin{frame}{Reweighting and translation give the same law identity}
With \(Z=e^{-\theta X-\theta^2/2}\),
\[
\boxed{
\E^\Q[f(X)]
=\E^\Pp[Zf(X)]
=\E^\Pp[f(X-\theta)].}
\]
\begin{pausebox}
Which equality is the definition of \(\Q\), and which equality uses the
Gaussian law?
\end{pausebox}
\pause
The first is change of measure; the second is a distributional consequence.
\end{frame}

\begin{frame}{A shifted tail needs a likelihood correction}
Recall \(p=\Pp(X>5)\approx2.87\times10^{-7}\). Sample
\(X_1,\ldots,X_n\) under \(\Q\), where \(X\sim\Nlaw(5,1)\).
\begin{columns}[T,onlytextwidth]
\column{0.48\textwidth}
\begin{alertblock}{Shift the mean only}
\[
\widehat p_{\mathrm{shift}}
=\frac1n\sum_{i=1}^n\1_{\{X_i>5\}}
\longrightarrow \Q(X>5)=\frac12.
\]
This estimates the tail under \(\Q\), not the original tail.
\end{alertblock}
\column{0.48\textwidth}
\begin{exampleblock}{Restore the original weights}
\[
w(x)=\frac{\dd\Pp}{\dd\Q}(x)
=e^{-5x+25/2},
\]
\[
\widehat p_{\mathrm{IS}}
=\frac1n\sum_{i=1}^n w(X_i)\1_{\{X_i>5\}}.
\]
\end{exampleblock}
\end{columns}
\medskip
\centering\(\boxed{\E^\Q[w(X)\1_{\{X>5\}}]=\Pp(X>5)=p.}\)
\end{frame}



\section{Why Change the Probability Measure?}

\begin{frame}{The payoff stays the same while scenario weights change}
For a terminal payoff \(H\), valuation depends on its possible values and on
the probabilities assigned to the scenarios.
\begin{exampleblock}{Change of viewpoint}
A change of measure does not modify the state, the stock path, or \(H(\omega)\).
It changes how strongly each state is weighted.
\end{exampleblock}
\begin{alertblock}{What Girsanov will do}
The drift may change. The diffusion coefficient will not.
\end{alertblock}
\end{frame}

\begin{frame}{A pricing measure is not a forecast}
\begin{columns}[T,onlytextwidth]
\column{0.48\textwidth}
\begin{modelbox}
Under \(\Pp\), dynamics describe statistical behavior:
\[
\frac{\dd S_t}{S_t}=\mu_t\dd t+\sigma_t\dd W_t^\Pp.
\]
\end{modelbox}
\column{0.48\textwidth}
\begin{marketbox}
Under \(\Q\), discounted traded prices should be martingales:
\[
\frac{\dd S_t}{S_t}=r_t\dd t+\sigma_t\dd W_t^\Q.
\]
\end{marketbox}
\end{columns}
\end{frame}

\begin{frame}{Two practical reasons to change the measure}
\begin{columns}[T,onlytextwidth]
\column{0.48\textwidth}
\begin{marketbox}
\textbf{Option pricing.} Reweight paths so discounted traded prices have
zero drift. Then a price is a discounted expectation under \(\Q\).
\end{marketbox}
\column{0.48\textwidth}
\begin{exampleblock}{Rare-event simulation}
Reweight paths toward a rare loss (or barrier hit), estimate it more often,
then correct the estimate with the likelihood ratio \(\dd\Pp/\dd\Q\).
\end{exampleblock}
\end{columns}
\medskip
\begin{block}{Same paths, different probabilities}
The underlying scenarios and payoff do not change; only their weights do.
The new measure is chosen to make the calculation or simulation easier.
\end{block}
\end{frame}

\begin{frame}{A measure change is more than changing one variable's law}
\begin{columns}[T,onlytextwidth]
\column{0.48\textwidth}
\begin{exampleblock}{Tilt one variable}
For \(X\), a density proportional to \(e^{\theta X}\) changes the
marginal distribution of \(X\). It says nothing by itself about how other
variables or the rest of the path are reweighted.
\end{exampleblock}
\column{0.48\textwidth}
\begin{block}{Change the measure}
Choose a positive random weight \(Z_T\) with \(\E^\Pp[Z_T]=1\), and set
\(\dd\Q=Z_T\dd\Pp\). Every event and every payoff is reweighted:
\[
\Q(A)=\E^\Pp[Z_T\1_A],\qquad
\E^\Q[H]=\E^\Pp[Z_TH].
\]
\end{block}
\end{columns}
\medskip
An exponential tilt of \(X\) \emph{can} define a measure change; Girsanov's
density extends that idea to entire Brownian paths.
\end{frame}

\begin{frame}{Discounting reveals the drift to remove}
Suppose \(B_t=e^{rt}\) and
\[
\frac{\dd S_t}{S_t}=\mu\dd t+\sigma\dd W_t.
\]
\begin{pausebox}
Compute the dynamics of \(\bar S_t=S_t/B_t\). Which term must disappear?
\end{pausebox}
\pause
\[
\boxed{\frac{\dd\bar S_t}{\bar S_t}
=(\mu-r)\dd t+\sigma\dd W_t.}
\]
The ratio \((\mu-r)/\sigma\) is excess drift measured in volatility units.
\end{frame}



\section{From Gaussian Tilting to Brownian Paths}

\begin{frame}{A deterministic integrand produces one Gaussian variable}
Let \(h\in L^2([0,T])\) be deterministic and define
\[
X=\int_0^T h_s\dd W_s
\sim\Nlaw\!\left(0,\int_0^T h_s^2\dd s\right).
\]
Set \(v=\int_0^T h_s^2\dd s\). The Gaussian MGF says
\(\E[e^X]=e^{v/2}\).
\begin{pausebox}
Find a unit-mean exponential weight built from \(X\), and write it as a
density for reweighting Brownian paths.
\end{pausebox}
\pause
The normalizing correction is \(-v/2\), so the path density is
\[
\boxed{
Z_T=\exp\!\left(
\int_0^T h_s\dd W_s-\frac12\int_0^T h_s^2\dd s
\right).}
\]
\end{frame}


\begin{frame}{Cameron--Martin absorbs deterministic drift}
Define \(\Q\) by \(\dd\Q=Z_T\dd\Pp\). Then
\[
\boxed{
B_t=W_t-\int_0^th_s\dd s
\quad\text{is Brownian under }\Q.}
\]
Equivalently,
\[
\dd W_t=\dd B_t+h_t\dd t.
\]
\begin{block}{Pathwise interpretation}
The density rewards paths aligned with \(h\), so \(W\) acquires drift
\(\int h\) under \(\Q\).
\end{block}
\end{frame}

\begin{frame}{A constant drift checks the theorem against Gaussian tilting}
Let \(h_t\equiv\theta\). Then
\[
Z_T=e^{\theta W_T-\theta^2T/2},
\qquad B_t=W_t-\theta t.
\]
\begin{pausebox}
What is the law of \(W_T\) under \(\Q\)?
\end{pausebox}
\pause
Since \(W_T=B_T+\theta T\),
\[
\boxed{W_T\sim\Nlaw(\theta T,T)\text{ under }\Q.}
\]
\end{frame}

\begin{frame}{Random adapted drift requires a new argument}
For deterministic \(h\), \(\int h\dd W\) is Gaussian, so one mgf
normalizes the density.
\begin{alertblock}{What changes}
If \(h\) is replaced by a random adapted process \(\phi\), the stochastic
integral need not be Gaussian. We must use stochastic calculus to control
\(\E[Z_T]\).
\end{alertblock}
\end{frame}



\section{The Stochastic Exponential}

\begin{frame}{Adapted drift leads to an exponential process}
Let \(\phi\) be progressively measurable and suppose
\[
\int_0^T\phi_s^2\dd s<\infty\quad\text{a.s.}
\]
Define
\[
Z_t=\exp\!\left(
\int_0^t\phi_s\dd W_s-\frac12\int_0^t\phi_s^2\dd s
\right).
\]
\begin{pausebox}
Compute \(dZ_t\) before calculating.
\end{pausebox}
\end{frame}

\begin{frame}{It\^o's formula explains the compensation}
Set
\[
Y_t=\int_0^t\phi_s\dd W_s-\frac12\int_0^t\phi_s^2\dd s.
\]
Then
\[
\dd Y_t=\phi_t\dd W_t-\frac12\phi_t^2\dd t,
\qquad
\dd[Y]_t=\phi_t^2\dd t.
\]
Applying It\^o to \(Z_t=e^{Y_t}\),
\[
\dd Z_t=Z_t\dd Y_t+\frac12Z_t\dd[Y]_t.
\]
\end{frame}

\begin{frame}{Quadratic variation cancels the finite-variation drift}
\[
\begin{aligned}
\dd Z_t
&=Z_t\left(\phi_t\dd W_t-\frac12\phi_t^2\dd t\right)
+\frac12Z_t\phi_t^2\dd t\\
&=\boxed{Z_t\phi_t\dd W_t}.
\end{aligned}
\]
\begin{pausebox}
Does zero drift already prove that \(Z\) is a true martingale?
\end{pausebox}
\pause
No. It first gives a local martingale; integrability is still required.
\end{frame}


\begin{frame}{Novikov is a usable sufficient condition}
\begin{block}{Novikov condition}
If
\[
\boxed{
\E^\Pp\!\left[
\exp\!\left(\frac12\int_0^T\phi_s^2\dd s\right)
\right]<\infty,}
\]
then \(Z\) is a true martingale and \(\E^\Pp[Z_T]=1\).
\end{block}
We use the condition; its proof is beyond today's scope.
\end{frame}

\begin{frame}{Bounded drift makes Novikov immediate}
If \(|\phi_t|\le C\), then
\[
\int_0^T\phi_s^2\dd s\le C^2T
\]
and
\[
\E\!\left[e^{\frac12\int_0^T\phi_s^2\dd s}\right]
\le e^{C^2T/2}<\infty.
\]
\begin{pausebox}
Does finite energy almost surely imply Novikov automatically?
\end{pausebox}
\pause
No. Almost-sure finiteness does not control an exponential expectation.
\end{frame}

\begin{frame}{Diagnose three candidate densities}
\small
\[
\begin{array}{c|c}
\text{candidate} & \text{information}\\ \hline
e^{\theta W_T-\theta^2T/2} & \theta\in\R\\
e^{\int_0^T\phi_s\dd W_s-\frac12\int_0^T\phi_s^2\dd s}
& |\phi_s|\le C\\
e^{\int_0^T\psi_s\dd W_s-\frac12\int_0^T\psi_s^2\dd s}
& \int_0^T\psi_s^2\dd s<\infty\text{ a.s.}
\end{array}
\]
\begin{pausebox}
Which can immediately define a new probability measure? Justify.
\end{pausebox}
\pause
The first is Gaussian, the second satisfies Novikov, and the third needs
more information.
\end{frame}

\section{Girsanov Moves Drift into the Measure}

\begin{frame}{Girsanov identifies the new Brownian motion}
Assume \(Z\) is a true martingale and
\[
\frac{\dd\Q}{\dd\Pp}\bigg|_{\cF_T}=Z_T,
\qquad
Z_t=e^{\int_0^t\phi_s\dd W_s-\frac12\int_0^t\phi_s^2\dd s}.
\]
\begin{block}{Girsanov theorem}
Under \(\Q\),
\[
\boxed{B_t=W_t-\int_0^t\phi_s\dd s}
\]
is Brownian motion.
\end{block}
\end{frame}

\begin{frame}{Write the differential substitution before using the theorem}
The measure is defined by the full density process
\[
Z_t=\exp\!\left(\int_0^t\phi_s\dd W_s
-\frac12\int_0^t\phi_s^2\dd s\right),
\qquad \frac{\dd\Q}{\dd\Pp}\bigg|_{\cF_T}=Z_T.
\]
Under \(\Q\),
\[
B_t=W_t-\int_0^t\phi_s\dd s
\quad\Longleftrightarrow\quad
\boxed{\dd W_t=\dd B_t+\phi_t\dd t.}
\]
\begin{simplificationbox}
The density sign fixes the Brownian shift. State the theorem and solve for
\(\dd W_t\) before substituting into an SDE.
\end{simplificationbox}
\end{frame}

\begin{frame}{A generic It\^o process keeps its volatility}
Under \(\Pp\), suppose
\[
\dd X_t=b_t\dd t+a_t\dd W_t.
\]
Use the density
\[
Z_t=\exp\!\left(\int_0^t\phi_s\dd W_s
-\frac12\int_0^t\phi_s^2\dd s\right),
\qquad \frac{\dd\Q}{\dd\Pp}\bigg|_{\cF_T}=Z_T.
\]
Then under \(\Q\),
\[
\boxed{
\dd X_t=(b_t+a_t\phi_t)\dd t+a_t\dd B_t.}
\]
\begin{exampleblock}{Invariant}
Only the drift changes. The diffusion coefficient remains \(a_t\).
\end{exampleblock}
\end{frame}

\begin{frame}{Choose the shift to remove a target drift}
Let \(a_t\ne0\) in
\[
\dd X_t=b_t\dd t+a_t\dd W_t.
\]
\begin{pausebox}
Choose \(\phi_t\) so that \(X\) has zero drift under \(\Q\).
\end{pausebox}
\pause
\[
\boxed{\phi_t=-\frac{b_t}{a_t}.}
\]
The corresponding density is
\[
Z_t=\exp\!\left(-\int_0^t\frac{b_s}{a_s}\dd W_s
-\frac12\int_0^t\frac{b_s^2}{a_s^2}\dd s\right),
\qquad \frac{\dd\Q}{\dd\Pp}\bigg|_{\cF_T}=Z_T.
\]
This is a formal candidate until \(Z\) is shown to be a true martingale.
\end{frame}

\begin{frame}{The theorem becomes useful through three translations}
Girsanov has identified a new Brownian motion. We now use it to answer
concrete questions.
\begin{enumerate}
  \item Read the Brownian shift from the density
  \(Z_t=\exp\!\bigl(\int_0^t\phi_s\dd W_s
  -\frac12\int_0^t\phi_s^2\dd s\bigr)\).
  \item Substitute \(\dd W\) and write the full SDE under \(\Q\).
  \item Use the new SDE to identify a law or compute an expectation.
\end{enumerate}
\begin{block}{The progression}
We start with \(W\), then an arithmetic diffusion, an adapted drift, and
finally a geometric asset price.
\end{block}
\end{frame}

\begin{frame}{A constant density moves the Brownian mean}
Define
\[
Z_T=e^{-\lambda W_T-\lambda^2T/2}.
\]
\begin{pausebox}
Find the \(\Q\)-Brownian motion and the law of \(W_T\) under \(\Q\).
\end{pausebox}
\pause
Here \(\phi=-\lambda\), so
\[
W_t^\Q=W_t+\lambda t,\qquad W_T=W_T^\Q-\lambda T,
\]
and therefore
\[
\boxed{W_T\sim\Nlaw(-\lambda T,T)\text{ under }\Q.}
\]
\end{frame}

\begin{frame}{The same shift removes drift from an arithmetic diffusion}
Under \(\Pp\), let
\[
\dd X_t=\mu\,\dd t+\sigma\,\dd W_t,
\qquad \sigma>0,
\]
and set \(\lambda=\mu/\sigma\).
\[
Z_T=e^{-\lambda W_T-\lambda^2T/2},
\qquad
W_t^\Q=W_t+\lambda t.
\]
Since \(\dd W_t=\dd W_t^\Q-\lambda\dd t\),
\[
\boxed{\dd X_t=\sigma\,\dd W_t^\Q.}
\]
If \(X_0=x\), then under \(\Q\),
\[
\boxed{X_t\sim\Nlaw(x,\sigma^2t).}
\]
\end{frame}

\begin{frame}{A change of measure can target any constant drift}
Suppose again
\[
\dd X_t=\mu\,\dd t+\sigma\,\dd W_t,
\]
but now we want drift \(\nu\), not necessarily zero.
\begin{pausebox}
Which Brownian shift produces \(\dd X_t=\nu\dd t+\sigma\dd W_t^\Q\)?
\end{pausebox}
\pause
Choose
\[
\lambda=\frac{\mu-\nu}{\sigma},
\qquad
Z_t=e^{-\lambda W_t-\lambda^2t/2},
\qquad
\frac{\dd\Q}{\dd\Pp}\bigg|_{\cF_T}=Z_T.
\]
Then \(W_t^\Q=W_t+\lambda t\) and
\[
\boxed{\dd X_t=\nu\dd t+\sigma\dd W_t^\Q.}
\]
\end{frame}

\begin{frame}{A bounded adapted drift can also be absorbed}
Under \(\Pp\), consider
\[
\dd X_t=\sin(W_t)\dd t+\dd W_t,\qquad X_0=0.
\]
Take \(\phi_t=-\sin(W_t)\) and define \(\Q\) by
\[
Z_t=\exp\!\left(
-\int_0^t\sin(W_s)\dd W_s
-\frac12\int_0^t\sin^2(W_s)\dd s
\right),\qquad
\frac{\dd\Q}{\dd\Pp}\bigg|_{\cF_T}=Z_T.
\]
Because \(|\phi_t|\le1\), Novikov holds. Under \(\Q\),
\[
B_t=W_t+\int_0^t\sin(W_s)\dd s
\]
is Brownian, and substitution gives
\[
\boxed{\dd X_t=\dd B_t.}
\]
\end{frame}

\begin{frame}{A geometric diffusion prepares the pricing measure}
Under \(\Pp\), suppose
\[
\frac{\dd S_t}{S_t}=\mu\dd t+\sigma\dd W_t,
\qquad \sigma>0.
\]
To target a new drift \(\nu\), choose
\[
\lambda=\frac{\mu-\nu}{\sigma},
\qquad
Z_t=e^{-\lambda W_t-\lambda^2t/2},
\qquad
\frac{\dd\Q}{\dd\Pp}\bigg|_{\cF_T}=Z_T.
\]
Then \(W_t^\Q=W_t+\lambda t\) and
\[
\boxed{\frac{\dd S_t}{S_t}
=\nu\dd t+\sigma\dd W_t^\Q.}
\]
\begin{marketbox}
For pricing with the bank account as numeraire, the required target is
\(\nu=r\). This is the next step.
\end{marketbox}
\end{frame}

\begin{frame}{New dynamics give laws; the density gives conditioning}
Once the SDE is written under \(\Q\), we can work directly with its new law.
To translate a conditional expectation back to \(\Pp\), use
\[
\boxed{
\E^\Q[Y\mid\cF_s]
=\frac1{Z_s}\E^\Pp[Z_TY\mid\cF_s],
\qquad Z_s=\E^\Pp[Z_T\mid\cF_s].}
\]
For the geometric example with constant \(\nu,\sigma\),
\[
\E^\Q[S_T\mid\cF_s]
=S_se^{\nu(T-s)}
=\frac1{Z_s}\E^\Pp[Z_TS_T\mid\cF_s].
\]
The new SDE gives the middle expression; the density translates it back to
the historical measure.
\end{frame}

\begin{frame}{The examples reveal one reusable workflow}
\begin{enumerate}
  \item \textbf{Target:} identify the drift wanted under the new measure.
  \item \textbf{Density:} write \(Z_T\) with the correct sign and compensator.
  \item \textbf{Validity:} prove or assume that \(Z\) is a true martingale.
  \item \textbf{Substitution:} identify \(W^\Q\), solve for \(\dd W\), and
  write the SDE under \(\Q\).
  \item \textbf{Use:} compute the new law or a conditional expectation.
\end{enumerate}
\begin{pausebox}
Which step remains unresolved if Novikov is unavailable?
\end{pausebox}
\pause
Only validity; the algebraic drift candidate still exists.
\end{frame}

\section{The Risk-Neutral Measure in Black--Scholes}

\begin{frame}{Discounting exposes the drift that must be removed}
Under \(\Pp\),
\[
\frac{\dd S_t}{S_t}=\mu_t\dd t+\sigma_t\dd W_t,
\qquad
B_t^{(0)}=e^{\int_0^tr_s\dd s},
\qquad D_t=(B_t^{(0)})^{-1}.
\]
Since \(D\) has finite variation,
\[
\boxed{
\frac{\dd(D_tS_t)}{D_tS_t}
=(\mu_t-r_t)\dd t+\sigma_t\dd W_t.}
\]
\end{frame}

\begin{frame}{The market price of risk is the required shift}
Define
\[
\lambda_t=\frac{\mu_t-r_t}{\sigma_t}.
\]
Then
\[
\frac{\dd(D_tS_t)}{D_tS_t}
=\lambda_t\sigma_t\dd t+\sigma_t\dd W_t.
\]
\begin{pausebox}
Should the density contain \(+\int\lambda\dd W\) or
\(-\int\lambda\dd W\) to eliminate the drift?
\end{pausebox}
\pause
It must use the negative sign.
\end{frame}

\begin{frame}{The pricing density has the negative market price of risk}
\[
\boxed{
Z_t=\exp\!\left(
-\int_0^t\lambda_s\dd W_s
-\frac12\int_0^t\lambda_s^2\dd s
\right).}
\]
Assume \(Z\) is a true martingale and define
\[
\frac{\dd\Q}{\dd\Pp}\bigg|_{\cF_T}=Z_T.
\]
Girsanov gives
\[
\boxed{W_t^\Q=W_t+\int_0^t\lambda_s\dd s.}
\]
\end{frame}

\begin{frame}{Substitution replaces \(\mu\) by \(r\)}
\[
\dd W_t=\dd W_t^\Q-\lambda_t\dd t.
\]
Therefore
\[
\begin{aligned}
\frac{\dd S_t}{S_t}
&=\mu_t\dd t+\sigma_t(\dd W_t^\Q-\lambda_t\dd t)\\
&=(\mu_t-\sigma_t\lambda_t)\dd t+\sigma_t\dd W_t^\Q\\
&=\boxed{r_t\dd t+\sigma_t\dd W_t^\Q}.
\end{aligned}
\]
\end{frame}

\begin{frame}{Discounted stock prices lose their drift under \(\Q\)}
Under \(\Q\),
\[
\frac{\dd S_t}{S_t}=r_t\dd t+\sigma_t\dd W_t^\Q.
\]
Thus
\[
\boxed{
\dd(D_tS_t)=\sigma_tD_tS_t\dd W_t^\Q.}
\]
\begin{block}{Conclusion}
The discounted stock is a \(\Q\)-local martingale and, under suitable
integrability, a true martingale.
\end{block}
\end{frame}

\begin{frame}{Risk-neutral does not describe investor preferences}
\begin{columns}[T,onlytextwidth]
\column{0.48\textwidth}
\begin{block}{Correct meaning}
\(\Q\) is an equivalent measure under which discounted traded prices are
martingales.
\end{block}
\column{0.48\textwidth}
\begin{alertblock}{Incorrect meaning}
\(\Q\) does not claim that real investors expect every asset to grow at
the risk-free rate.
\end{alertblock}
\end{columns}
\medskip
It is a pricing representation/ a tool enforced/given by no arbitrage.
\end{frame}

\begin{frame}{Split the call payoff into two probability questions}
Under the risk-neutral measure,
\[
S_T=s\exp\!\left((r-\tfrac12\sigma^2)\tau+\sigma W_\tau^\Q\right),
\qquad C=e^{-r\tau}\E^\Q[(S_T-K)^+].
\]
Write the payoff as \((S_T-K)\1_{\{S_T>K\}}\). Then
\[
C=e^{-r\tau}\E^\Q[S_T\1_{\{S_T>K\}}]
  -Ke^{-r\tau}\Q(S_T>K).
\]
\begin{pausebox}
Using the lognormal formula, show that
\(\Q(S_T>K)=N(d_-)\), where
\(d_-=\frac{\log(s/K)+(r-\frac12\sigma^2)\tau}
{\sigma\sqrt\tau}\).
\end{pausebox}
\end{frame}

\begin{frame}{Girsanov interprets the stock-weighted tail}
Define the stock-numeraire measure by
\[
\left.\frac{\dd\Q^{S}}{\dd\Q}\right|_{\cF_T}
=\frac{e^{-r\tau}S_T}{s}
=\exp\!\left(\sigma W_\tau^\Q-\tfrac12\sigma^2\tau\right).
\]
This is a unit-mean density. By Girsanov,
\[
W_u^{S}=W_u^\Q-\sigma u
\quad\text{is Brownian under }\Q^{S}.
\]
Thus the first payoff contribution becomes
\[
e^{-r\tau}\E^\Q[S_T\1_{\{S_T>K\}}]
=s\,\Q^{S}(S_T>K).
\]
Under \(\Q^S\), the same event \(S_T>K\) has probability \(N(d_+)\), with
\[
d_+=\frac{\log(s/K)+(r+\frac12\sigma^2)\tau}{\sigma\sqrt\tau}.
\]
\end{frame}

\begin{frame}{The call formula is a risk-neutral Gaussian expectation}
For \(\tau=T-t\),
\[
C(t,s)=s\,N(d_+)-Ke^{-r\tau}N(d_-),
\]
\[
d_\pm
=\frac{\log(s/K)+(r\pm\frac12\sigma^2)\tau}
{\sigma\sqrt{\tau}},
\qquad d_-=d_+-\sigma\sqrt{\tau}.
\]
\begin{block}{Two measures, same exercise event}
\(N(d_-)=\Q(S_T>K)\) under the money-market measure, while
\(N(d_+)=\Q^S(S_T>K)\) under the stock-numeraire measure. Thus
\(sN(d_+)\) is the stock-weighted tail, and \(Ke^{-r\tau}N(d_-)\) is
the discounted strike tail.
\end{block}
\end{frame}

\begin{frame}{Delta turns a pricing formula into a trading rule}
\[
\boxed{\Delta_t=C_s(t,S_t)=N(d_+(t,S_t)).}
\]
\begin{pausebox}
As maturity approaches, what should delta do when the call is deeply
in the money? Deeply out of the money? Near the strike?
\end{pausebox}
\pause
\[
\Delta_t\longrightarrow
\begin{cases}
1,&S_t>K,\\
0,&S_t<K,
\end{cases}
\]
while the transition near \(K\) becomes increasingly sharp.
\end{frame}

\begin{frame}{Complete the Black--Scholes sign drill}
Under \(\Pp\),
\[
\frac{\dd S_t}{S_t}=\mu\dd t+\sigma\dd W_t^\Pp.
\]
\begin{pausebox}
Write, in order:
\begin{enumerate}
  \item the market price of risk;
  \item the density process;
  \item the \(\Q\)-Brownian motion;
  \item the \(\Q\)-dynamics of \(S\);
  \item the dynamics of \(e^{-rt}S_t\).
\end{enumerate}
\end{pausebox}
\end{frame}

\end{document}
